\scanpage{14}
\mspage{2)}
Mais on ne peut pas dire \inserted{en général} que $\dot T$ soit la frontière
de $T$ dans $X$, i.e. (7) devient faux si $\lambda=1$
% P02: crossed group after T; the lambda reading and lowered attachment remain uncertain.
(prendre $X=T\deleted{\unsure{_\lambda}}$, \unsure{pour|avec} $\dot T\neq0$).
% P02: isolated vertical ink point below faux, to the right of the parenthesis; linguistic status unclear.
\figure{\begin{tikzpicture}[x=1cm,y=1cm]
\path[use as bounding box] (0,0) rectangle (10,0.2);
\fill (6.2,0.1) ellipse (0.025 and 0.06);
\end{tikzpicture}}

% P03: the isolated raised capital is not expanded into a word; adjacent slanted ink belongs to the marginal note.
\deleted{\unsure{On dira}} \inserted{\unsure{É|E}} \unsure{enfin de} dire qu'une partie fermée $\deleted{Z}Y$ d'un
espace topologique $X$ est une ``\uline{partie bord}''
s'il existe un \uline{voisinage}-\deleted{\uline{collier}}
% P05: the initial group is overworked; no independent earlier word can be separated from it.
\inserted{\uline{\unsure{c}ylindrique}} de $\deleted{Z}Y$ dans $X$,
i.e. un voisinage \deleted{tubulaire} cônique $T$ tel que
$\dot T\longrightarrow Y$ soit un homéom --- ou ce qui revient
au même, \deleted{\unsure{i.e.} tel que l'on} s'il existe un
voisinage $\deleted{\unsure{de}}T$ de $Y$ dans $X$ et un homéom
\[
(9)\qquad T\simeq Y\times\mathbf{I}
\]
rendant commutatif le diagramme
\diagram{\[
(10)\qquad
\begin{tikzcd}
Y\arrow[r,hook]\arrow[dr,"{y\mapsto(y,0)}"']&T\arrow[d,phantom,"{\wr\!\vert}" description]\\
&Y\times\mathbf{I}
\end{tikzcd}
\]}
\margin{définir la structure de voisinage cylindrique
comme un homéom.
$E\simeq Y\times[0,1]$
compatible avec
% P10: the upper connector has no secure arrowhead; the target is literally [0,1].
\diagram{\[
\begin{tikzcd}[column sep=small,row sep=small]
Y&\unsure{-|\longrightarrow}&[0,1]\\
y&\longmapsto&(y,0)
\end{tikzcd}
\]}}

Par exemple, si $T$ est un voisinage \deleted{tubulaire} cônique
\inserted{(absolu)} de \deleted{\unsure{tout}} $Y$, alors $\dot T$ est un sous-espace bord.
\unsure{donc} $\dot T_\lambda$ est un sous-espace bord de $T_\lambda$ ;
$\deleted{\unsure{\mathrm{Ep}}}\dot T,(1-\unsure{\varepsilon|\delta})\dot T$
et $\dot T\cup(1-\varepsilon)\dot T$ sont des sous-espaces
bords de $\bar{\mathrm{E}}\mathrm{p}_\varepsilon(T)$. \deleted{Dans}
Si $X$ est un espace
\uline{compact}, et $(Y_i)_{i\in I}$ une partie finie de
parties fermées disjointes de $X$, alors
% P14: the lower union index is J-shaped; the I alternative and the contrast with the earlier I remain explicit.
$Y=\bigcup_{i\in\unsure{J|I}}Y_i$ est un sous-espace bord ssi les $Y_i$ le sont
\margin{13}
